\documentclass[11pt]{article}
\usepackage{amsmath,amssymb}
\usepackage[hidelinks]{hyperref}
\title{A pending send and a spent send at a live repair cut}
\author{}
\date{}
\begin{document}
\maketitle
\paragraph{Scope.}
Protect one old warehouse completion report through a replacement of its
designated radio supplier. Pinned MRWS completion and pinned SubT dispatch
are coupled with deterministic path/pose seams, one optional datagram, and
one serialized rejoin boundary. The appointment, right to remove the
optional entry, ready backup and deadline are authored. The old body
comes from the native \texttt{ORDERCOMPLETE}; the adapter checks its
actual receipt. This is not exclusive-custody proof, field evidence,
or I1/R2 acceptance.

\paragraph{Dependency discovered at the effect source.}
The optional datagram goes from \(Q\) to Side and does not occupy Base
receive capacity. If a pose failure blocks dispatch, it remains queued and
there is no debit. If it leaves the deque but its path is erased, there is
no receive callback either, even though \(Q\)'s rolling debit is live.
Neither condition affects incumbent \(P\). Both can matter after the
appointment changes to \(Q\). A second failed pose update occurs while the
report is outstanding. Late recognition selects already authorized
fallback \(R\). Independent traffic at rejoin uses disjoint resources.

\paragraph{Exact cut.}
For the rejoin dispatch, let \(C\) be capacity, \(L\) the actual unexpired
charge to \(Q\), and \(K\in\{0,1\}\) indicate the pending optional send
kept ahead of the old report. Given viable paths and no other relevant
traffic, service by \(Q\) occurs iff \(L+60K+173\le C\).
At \(173\le C<233\), a pending send must be removed to use \(Q\); an already
charged send cannot be removed and requires another lawful sender or expiry
by deadline. The native charge remains live at exactly one second and is
gone at the next quarter-second tested. If the policy must protect the old
packet then avoid the backup when possible, absent callback alone cannot
attain this profile: the observation class has both histories. Conservative
use of the backup can still protect the report. At \(C\ge233\) both histories
permit \(Q\) without removal; below 173 no sender can send the report.
The optional probe carries no old obligation and is assumed withdrawable.

\paragraph{Discovery cost and control boundary.}
The adapter consults the actual deque and ledger, not a stipulated future
truth table, at a serialized cut. The identity scan inspects \(q\) queue
entries, the debit check scans a ledger of length \(b\), and removal may
scan and compact the deque: \(O(q+b)\) for bounded packet lengths. This is
not a bound for incremental index maintenance. Such an index would have to
cover enqueue, pop, charge, expiry and migration and remain coherent with
dispatch. If the optional send moves to the ledger after inspection but
before consumption, an unserialized inspection-plus-removal need not retain
the allowance. Recipient-side revocation checks, as in the
\href{https://www.usenix.org/legacy/event/osdi06/tech/full_papers/burrows/burrows.pdf}{Chubby sequencer},
provide a relevant prior control boundary, not a way to erase an effect
that occurs before the recipient. No stock API is claimed for the adapter's
queue surgery.

\paragraph{Comparison and residual.}
The same-information direct formula matches all 288 runs. In the 72 runs
per arm with \(C\ge173\), qualification gives 72 on-time receipts with 22
backup uses; using the backup uniformly when \(Q\) is selected at tight
capacity gives 72 with 30. Keeping the queue and omitting budget-based
substitution gives 64 and 18. Four qualified runs remove the named probe.
The improvement is in the usable information and action interface, not over
an equally informed scheduler. The native broker's all-member pose gate
still applies, so unrelated work proceeds at rejoin, not during that
blackout. Actual appointment provenance and recurrent maintenance of a
coherent observer, and an independent mission, are open.
\end{document}
